\chapter{V2 : le rayon de bord d'attaque comme variable}\label{ch:v2}
\iterbanner{V2}

\begin{evolution}{V2}{pourquoi ce chapitre}En V1, la SER de face d'un profil NACA dépendait presque uniquement de l'épaisseur, par le rayon de bord d'attaque $r_{BA}\propto t^2$ (chapitre~\ref{ch:naca}). En V2, la géométrie devient la série NACA modifiée (cours aéro, chapitre V2), où l'indice $I$ fixe le rayon indépendamment de l'épaisseur : $r_{BA}=1{,}1019\,(tI/6)^2$. Le modèle SER lui-même ne change pas.\end{evolution}

\section{La SER de nez en fonction de $I$}
\begin{figure}[htbp]
\centering
% figure générée par tools/make_figures_v2.py : v2_nez
\begin{tikzpicture}
\begin{axis}[width=0.72\linewidth,height=5.8cm,xlabel={indice de rayon de bord d'attaque $I$},ylabel={$\sigma_{2D}$ de nez (dBm)},
  xmin=1.5,xmax=9,legend style={draw=none,font=\scriptsize},legend pos=south east]
\addplot[corange,very thick] coordinates {(1.5,-25.683) (1.75,-24.879) (2,-24.134) (2.25,-23.442) (2.5,-22.796) (2.75,-22.191) (3,-21.623) (3.25,-21.088) (3.5,-20.583) (3.75,-20.103) (4,-19.648) (4.25,-19.215) (4.5,-18.802) (4.75,-18.406) (5,-18.028) (5.25,-17.665) (5.5,-17.316) (5.75,-16.98) (6,-16.657) (6.25,-16.345) (6.5,-16.043) (6.75,-15.752) (7,-15.47) (7.25,-15.197) (7.5,-14.931) (7.75,-14.674) (8,-14.424) (8.25,-14.181) (8.5,-13.944) (8.75,-13.714) (9,-13.489)};\addlegendentry{PO, $t=8$~\%}
\addplot[corange,dashed,forget plot] coordinates {(1.5,-28.586) (1.75,-27.248) (2,-26.088) (2.25,-25.065) (2.5,-24.15) (2.75,-23.322) (3,-22.566) (3.25,-21.871) (3.5,-21.227) (3.75,-20.628) (4,-20.067) (4.25,-19.541) (4.5,-19.044) (4.75,-18.574) (5,-18.129) (5.25,-17.705) (5.5,-17.301) (5.75,-16.915) (6,-16.545) (6.25,-16.191) (6.5,-15.85) (6.75,-15.522) (7,-15.206) (7.25,-14.902) (7.5,-14.607) (7.75,-14.322) (8,-14.047) (8.25,-13.779) (8.5,-13.52) (8.75,-13.268) (9,-13.023)};
\addplot[cblue,very thick] coordinates {(1.5,-22.161) (1.75,-21.357) (2,-20.612) (2.25,-19.92) (2.5,-19.274) (2.75,-18.67) (3,-18.102) (3.25,-17.566) (3.5,-17.061) (3.75,-16.582) (4,-16.126) (4.25,-15.693) (4.5,-15.28) (4.75,-14.885) (5,-14.506) (5.25,-14.143) (5.5,-13.794) (5.75,-13.458) (6,-13.135) (6.25,-12.823) (6.5,-12.522) (6.75,-12.23) (7,-11.948) (7.25,-11.675) (7.5,-11.41) (7.75,-11.152) (8,-10.902) (8.25,-10.659) (8.5,-10.422) (8.75,-10.192) (9,-9.9676)};\addlegendentry{PO, $t=12$~\%}
\addplot[cblue,dashed,forget plot] coordinates {(1.5,-25.065) (1.75,-23.726) (2,-22.566) (2.25,-21.543) (2.5,-20.628) (2.75,-19.8) (3,-19.044) (3.25,-18.349) (3.5,-17.705) (3.75,-17.106) (4,-16.545) (4.25,-16.019) (4.5,-15.522) (4.75,-15.053) (5,-14.607) (5.25,-14.183) (5.5,-13.779) (5.75,-13.393) (6,-13.023) (6.25,-12.669) (6.5,-12.328) (6.75,-12) (7,-11.685) (7.25,-11.38) (7.5,-11.085) (7.75,-10.8) (8,-10.525) (8.25,-10.257) (8.5,-9.9981) (8.75,-9.7463) (9,-9.5016)};
\addplot[cgreen,very thick] coordinates {(1.5,-19.662) (1.75,-18.858) (2,-18.114) (2.25,-17.421) (2.5,-16.776) (2.75,-16.171) (3,-15.603) (3.25,-15.068) (3.5,-14.562) (3.75,-14.083) (4,-13.628) (4.25,-13.194) (4.5,-12.781) (4.75,-12.386) (5,-12.007) (5.25,-11.644) (5.5,-11.295) (5.75,-10.96) (6,-10.636) (6.25,-10.324) (6.5,-10.023) (6.75,-9.7315) (7,-9.4493) (7.25,-9.176) (7.5,-8.9108) (7.75,-8.6534) (8,-8.4034) (8.25,-8.1603) (8.5,-7.9237) (8.75,-7.6933) (9,-7.4688)};\addlegendentry{PO, $t=16$~\%}
\addplot[cgreen,dashed,forget plot] coordinates {(1.5,-22.566) (1.75,-21.227) (2,-20.067) (2.25,-19.044) (2.5,-18.129) (2.75,-17.301) (3,-16.545) (3.25,-15.85) (3.5,-15.206) (3.75,-14.607) (4,-14.047) (4.25,-13.52) (4.5,-13.023) (4.75,-12.554) (5,-12.108) (5.25,-11.685) (5.5,-11.28) (5.75,-10.894) (6,-10.525) (6.25,-10.17) (6.5,-9.8294) (6.75,-9.5016) (7,-9.1857) (7.25,-8.8809) (7.5,-8.5865) (7.75,-8.3017) (8,-8.0259) (8.25,-7.7586) (8.5,-7.4993) (8.75,-7.2475) (9,-7.0029)};
\end{axis}
\end{tikzpicture}

\caption{SER dans l'axe de profils symétriques modifiés, selon l'indice $I$ (10~GHz). Traits pleins : optique physique ; tirets : loi $\pi r_{BA}$. Entre $I=8$ et $I=2$, l'écho de nez baisse d'environ 10~dB à épaisseur fixée.}\label{fig:v2-nez}
\end{figure}

À épaisseur fixée, le rayon varie comme $I^2$ : passer de $I=6$ (NACA standard) à $I=2$ divise le rayon par 9. La loi du point spéculaire prédit $-9{,}5$~dB ; l'optique physique donne $-7{,}5$~dB à toutes les épaisseurs (par exemple de $-13{,}1$ à $-20{,}6$~dBm à $t=12$~\%). Pour la première fois, la géométrie offre un levier sur la SER \emph{sans toucher à l'épaisseur}.

\begin{figure}[htbp]
\centering
% figure générée par tools/make_figures_v2.py : v2_carte
\begin{tikzpicture}
\begin{axis}[width=0.5\linewidth,height=6cm,xlabel={indice $I$},ylabel={épaisseur $t$ (\%)},
  xtick={0,1,2,3},xticklabels={2,4,6,8},ytick={6,8,...,20},enlargelimits=false,axis on top,
  colormap={bl}{color=(white) color=(cblue!45) color=(cblue!85!black)},colorbar,
  colorbar style={ylabel={SER S1 (dBm)},width=2.5mm,tick label style={font=\scriptsize}},tick label style={font=\scriptsize}]
\addplot[matrix plot*,mesh/cols=4,point meta=explicit] coordinates {(0,6) [-26.108] (1,6) [-21.571] (2,6) [-18.655] (3,6) [-16.578] (0,8) [-23.454] (1,8) [-19.088] (2,8) [-16.222] (3,8) [-14.169] (0,10) [-21.462] (1,10) [-17.102] (2,10) [-14.306] (3,10) [-12.353] (0,12) [-19.794] (1,12) [-15.477] (2,12) [-12.764] (3,12) [-10.865] (0,14) [-18.364] (1,14) [-14.096] (2,14) [-11.470] (3,14) [-9.641] (0,16) [-17.106] (1,16) [-12.896] (2,16) [-10.362] (3,16) [-8.624] (0,18) [-16.010] (1,18) [-11.856] (2,18) [-9.396] (3,18) [-7.712] (0,20) [-15.023] (1,20) [-10.930] (2,20) [-8.512] (3,20) [-6.897]};
\end{axis}
\end{tikzpicture}

\caption{SER moyenne du secteur S1 sur la grille V2 (médiane sur la cambrure et sa position, à l'incidence de croisière). L'épaisseur et le rayon de bord d'attaque jouent désormais des rôles comparables (corrélations 0,71 et 0,67).}\label{fig:v2-carte}
\end{figure}

\section{Attention : les nez aigus sortent du régime optique}
Le prix de ce levier est un problème de validité. Pour $I=2$, à 10~GHz :
\begin{center}
\begin{tabular}{lcc}
\toprule
 & $t=8$~\% & $t=12$~\% \\
\midrule
$r_{BA}$ & 0,8~mm & 1,8~mm \\
$k\,r_{BA}$ & 0,16 & 0,37 \\
écart Mie / PO pour un cylindre (TM) & $+4{,}6$~dB & $+2{,}5$~dB \\
\bottomrule
\end{tabular}
\end{center}
Le nez est alors bien plus petit que la longueur d'onde : on n'est plus en régime optique, mais dans la zone de résonance, voire de Rayleigh. La solution exacte du cylindre montre qu'en polarisation TM, l'écho réel est \emph{plus fort} que $\pi a$, de 2,5 à 4,6~dB. Le gain de SER d'un nez aigu prédit par l'optique physique est donc probablement surestimé, peut-être de moitié en TM ; en TE, l'écart est de signe opposé.

\section{La réponse exacte : méthode des moments}
\begin{evolution}{V2}{section ajoutée après le premier calcul de la V2}Pour lever le doute, on a écrit une méthode des moments 2D en polarisation TM (\texttt{Code/ser/mom2d.py}), c'est-à-dire la résolution exacte de l'équation intégrale du champ électrique sur le contour. Elle reproduit la série de Mie du cylindre à 0,02~dB près, de $ka=0{,}5$ à $ka=31$, et ses résultats sur les profils sont convergés en maillage ($\lambda/15$, raffiné à $r_{BA}/4$ au nez) à 0,02~dB.\end{evolution}

\begin{figure}[htbp]
\centering
% figure générée par tools/make_figures_v2.py : v2_mom
\begin{tikzpicture}
\begin{axis}[width=0.72\linewidth,height=5.8cm,xlabel={indice de rayon de bord d'attaque $I$},ylabel={SER moyenne S1 (dBm)},
  xmin=1.5,xmax=8.5,legend style={draw=none,font=\scriptsize},legend pos=south east]
\addplot[corange,very thick,mark=*,mark size=1.4pt] coordinates {(2,-19.455) (3,-18.361) (4,-17.323) (6,-15.44) (8,-13.823)};\addlegendentry{MoM, $t=8$~\%}
\addplot[corange,dashed,mark=o,mark size=1.4pt,forget plot] coordinates {(2,-23.76) (3,-21.275) (4,-19.321) (6,-16.401) (8,-14.27)};
\addplot[cblue,very thick,mark=*,mark size=1.4pt] coordinates {(2,-17.679) (3,-16.191) (4,-14.846) (6,-12.585) (8,-10.775)};\addlegendentry{MoM, $t=12$~\%}
\addplot[cblue,dashed,mark=o,mark size=1.4pt,forget plot] coordinates {(2,-20.159) (3,-17.68) (4,-15.753) (6,-12.905) (8,-10.876)};
\addplot[cgreen,very thick,mark=*,mark size=1.4pt] coordinates {(2,-16.001) (3,-14.255) (4,-12.739) (6,-10.322) (8,-8.5093)};\addlegendentry{MoM, $t=16$~\%}
\addplot[cgreen,dashed,mark=o,mark size=1.4pt,forget plot] coordinates {(2,-17.545) (3,-15.083) (4,-13.187) (6,-10.441) (8,-8.5384)};
\end{axis}
\end{tikzpicture}

\caption{SER moyenne du secteur S1 de profils symétriques : méthode des moments (traits pleins) et optique physique (tirets). Les deux méthodes coïncident pour les nez ronds ; pour les nez aigus, l'optique physique sous-estime l'écho de 2,5 à 4,3~dB.}\label{fig:v2-mom}
\end{figure}

\paragraph{Ce que change la méthode exacte.}
\begin{itemize}
\item Le \textbf{gain d'un nez aigu existe}, mais il est plus faible : de $I=6$ à $I=2$, $-5{,}1$~dB à $t=12$~\% (au lieu de $-7{,}3$~dB) et $-4{,}0$~dB à $t=8$~\% (au lieu de $-7{,}4$~dB).
\item Sur les 1\,024 profils, l'optique physique sous-estime la SER de 1,05~dB en moyenne (écart-type 1,3~dB), mais le \textbf{classement} est presque identique : corrélation de rang 0,999.
\item Les \textbf{fronts de Pareto} sont robustes : 3 profils communs sur 4 dans S1, 9 sur 10 dans S2. Seul l'écart de SER le long du front diminue : de 1,8 à 0,8~dB dans S1, de 5,3 à 3,4~dB dans S2.
\end{itemize}

\begin{projet}
La méthode des moments devient la SER de référence de la V2 (\code{with\_objective(..., source="MoM")}) ; l'optique physique reste l'outil rapide. Le compromis traînée/SER de la V2 survit à la méthode exacte, avec un bénéfice radar plus modeste. Il reste à vérifier la polarisation TE et le passage en 3D, ce que le solveur du groupe SER, ou \texttt{bempp-cl} sur le STL, pourra faire.
\end{projet}

\exercices
\begin{exercice}
À quelle fréquence le nez d'un profil $t=12$~\%, $I=2$ (corde 1~m) atteindrait-il $k\,r_{BA}=3$, où l'optique physique est à 0,3~dB près ? Quelle est la longueur d'onde correspondante ?
\end{exercice}

\begin{pistes}
\begin{itemize}
\item \iterbadge{V2}~\textbf{Fait} : méthode des moments 2D en TM. Reste la polarisation TE (équation intégrale du champ magnétique), plus délicate.
\item Calculer la SER dans les deux polarisations et retenir la pire, ce qui est l'usage pour une menace de polarisation inconnue.
\end{itemize}
\end{pistes}
